% =============================================================================
% EINLEITUNG ZU BAND 5: SCHICHTEN, HILBERTRAUM-BRÜCKE UND JÜNGSTE KLÄRUNGEN
% (Bd. 5 von 5; Fortsetzung von Bd. 4)
% =============================================================================

\section*{Fünfter Band der Dokumentensammlung}

Dieser fünfte Band schließt die Sammlung von Einzeldokumenten zur T0-Theorie
/ FFGFT ab. Er versammelt jene Arbeiten, die im Korpus die methodische
Konsolidierung der Theorie tragen und in besonderem Maße ihren heutigen
Stand markieren.

\subsection*{Verhältnis zu Band 4}

Band 4 und Band 5 zusammen ersetzen das ursprünglich als Einzelband geplante
Erweiterungsbuch, das den Kindle-Umfang von 550 Seiten überschritten hätte.
Der Übergang von Band~4 nach Band~5 ist nicht arbiträr, sondern folgt einem
inhaltlichen Schnitt: bis Dok.~184 (p-bit) wird die Theorie noch wesentlich
auf Schicht~1 (kompaktes $T^4$, korrekturfreie Beschreibung in natürlichen
Einheiten) erarbeitet; ab Dok.~185 (Einbettungspreis) wird die
Schicht-1-/Schicht-2-Unterscheidung als ausdrückliches Beschreibungsraster
verwendet und bis zur epistemischen Selbstpositionierung in Dok.~262
durchgeführt.

\subsection*{Charakter dieses Bandes}

Band 5 ist dichter als die früheren Bände, weil seine Dokumente das
methodisch-grundlegende Gerüst der jüngsten Entwicklungsphase tragen:
\begin{itemize}
\item Die Hilbertraum-Brücke (Dok.~230--232), die FFGFT formal zur
      Quantenmechanik in Beziehung setzt und die operationale Äquivalenz
      mit interpretativer Divergenz festhält.
\item Die Schichten- und Skalenleiter-Dokumente (Dok.~241--253), die die
      Beschreibung in Schicht~1 (statisch, kompakt) und Schicht~2 (dynamisch,
      SI-projiziert) systematisch durcharbeiten.
\item Die jüngsten Klärungen (Dok.~255--262), die in der epistemischen
      Selbstpositionierung von Dok.~262 münden (Akzeptanz ohne Anschauung
      und Innensicht der dekompaktifizierten Zeit).
\end{itemize}

\subsection*{Gliederung in fünf Teile}

Die 37 Dokumente sind in numerischer Reihenfolge angeordnet und in fünf
thematische Teile gegliedert:

\begin{itemize}
\item \textbf{Teil I --- Späte Geometrie und Schichten-Vorbereitung}:
      Einbettungspreis, FFGFT-Photonik mit $\mathrm{LiNbO}_3$, Geometrische
      Grundlagen, Torus-Ableitung, Korrekturen-Register, Gemini-Dialog,
      Algebraische Kompositionsgrenzen, Nichtabschluss-Theorem.
\item \textbf{Teil II --- FFGFT-Feldtheorie und Operatoren}: Vollständige
      Feldtheorie-Darstellung, Rekursions-Operator $\mathcal{R}$, Fraktale
      Randbedingung, Narrativ, Dreieck-Matrix-Reduktion, Bewusstseins-Brücken,
      Korrekturen und Wicklungen.
\item \textbf{Teil III --- Hilbertraum-Brücke}: Hilbertraum-Übersetzung,
      Hilbertraum-Erweiterung, QG als Untermenge.
\item \textbf{Teil IV --- Schichten und Informationsformalismus}: Zwei
      Schichten, Skalenleiter, Informationsformalismus, UMF-Analyse,
      RA-Diagnostik, RSG-Vergleich, Information als Zustand und Prozess,
      Epistemologie, Falsifizierbarkeit, Information und SL,
      Vergleiche mit Vopson und Phillips, Xi-Zahlenraum-Analyse.
\item \textbf{Teil V --- Jüngste Klärungen und epistemische
      Selbstpositionierung}: Information aus Kinetik, Informationseinheit
      und Skala, Koide-Faktor 2/3, Koide-Kreuzterme, Statische natürliche
      Einheiten und SI-Projektion, sowie Akzeptanz ohne Anschauung und
      Innensicht der dekompaktifizierten Zeit (Dok.~262 als
      methodisch-grundlegender Abschluss).
\end{itemize}

\subsection*{Stand}

Sämtliche Dokumente dieses Bandes tragen den Stand Mai 2026.
